\chapter{Bộ công cụ ràng buộc đếm và giả Boolean}
\chapterlead{\emph{Cardinality constraint} (ràng buộc đếm) giới hạn số literal
đúng trong một tập; \emph{pseudo-Boolean constraint} (ràng buộc giả Boolean)
mở rộng mô hình này
bằng các trọng số nguyên. Hai lớp này xuất hiện trong ràng buộc gán, sức chứa,
chuỗi, đối xứng và hàm mục tiêu. Không có một bộ biểu diễn tốt nhất cho mọi miền;
cần chọn theo \(n\), ngưỡng \(k\), các đầu ra cần đọc và cách bộ giải thay cận.}

\index{cardinality constraint}\index{at-most-one}
\index{sequential counter}\index{totalizer}\index{sorting network}
\index{binary decision diagram}\index{pseudo-Boolean constraint}

\flowdiagram{AMO/EXO}{AMK/ALK}{Đúng \(K\)}
  {Giả Boolean}{Chọn bộ biểu diễn}

\section{Bản đồ phát triển của phép biểu diễn ràng buộc đếm}

Các phép biểu diễn ràng buộc đếm không phát triển theo một chuỗi trong đó bộ mã
hóa mới thay thế hoàn toàn bộ biểu diễn cũ. Mỗi họ giải quyết một kiểu đánh đổi:
\begin{enumerate}
  \item \emph{pairwise encoding} (biểu diễn từng cặp) và các phép phân rã AMO ưu tiên
        phát hiện xung đột trực tiếp;
  \item \emph{Sequential Counter} (bộ đếm tuần tự) hiện thực hóa (reification) tổng trên từng tiền tố;
  \item \emph{Totalizer} (cây tổng) hiện thực hóa (reification) tổng theo cấu trúc cây;
  \item \emph{Sorting Network / Selection Network} (mạng sắp xếp / lựa chọn) ưu tiên cấu
        trúc đều và độ sâu nhỏ;
  \item \emph{binary decision diagram} (biểu đồ quyết định nhị phân, BDD)
        hợp nhất các trạng thái tổng tương đương;
  \item Totalizer tổng quát và các biến thể PB xử lý trọng số;
  \item danh mục phương pháp hoặc bộ chọn đã học chọn phép biểu diễn theo đặc
        trưng của bài toán thử.
\end{enumerate}

Các khảo cứu về ràng buộc đếm nhấn mạnh rằng phải tách ba thuộc tính: kích
thước tiệm cận, mức nhất quán do lan truyền đơn vị duy trì và hiệu năng trong
một bộ giải cụ thể \parencite{bailleux2010survey}. Một Sequential Counter
\(O(nk)\) có thể nhỏ khi \(k\) bé; một mạng đếm có thể lớn hơn nhưng có độ sâu
tốt; Totalizer lại thuận lợi khi cận thay đổi. Không có quan hệ thứ tự tuyến
tính giữa ba lựa chọn đó.

\section{Các dạng ràng buộc}

Với \(X=\set{x_1,\ldots,x_n}\), ba tên tiếng Anh lần lượt là
\emph{at-most-\(k\)} (nhiều nhất \(k\), AMK), \emph{at-least-\(k\)}
(ít nhất \(k\), ALK) và \emph{exactly-\(k\)} (đúng \(k\), EXK):
\[
  \operatorname{AMK}(X,k):\quad \sum_{i=1}^{n}x_i\le k,
\]
\[
  \operatorname{ALK}(X,k):\quad \sum_{i=1}^{n}x_i\ge k,
\]
\[
  \operatorname{EXK}(X,k):\quad \sum_{i=1}^{n}x_i=k.
\]
\emph{At-most-one} (nhiều nhất một, AMO) là AMK với \(k=1\);
\emph{at-least-one} (ít nhất một, ALO) là một mệnh đề dài
\(x_1\lor\cdots\lor x_n\); \emph{exactly-one} (đúng một, EXO) là AMO kết hợp ALO.

Đối ngẫu hữu ích:
\[
  \sum_i x_i\ge k
  \iff
  \sum_i(1-x_i)\le n-k.
\]
Tuy nhiên, đối ngẫu có thể đổi ngưỡng từ nhỏ sang lớn. Vì vậy, lựa chọn giữa
Direct Encoding ALK và biểu diễn AMK trên các literal phủ định phụ thuộc vào
\(k\) và những đầu ra đã có.

\section{At-most-one (AMO): ràng buộc nhiều nhất một}

\subsection{Biểu diễn từng cặp}

Biểu diễn từng cặp sinh
\[
  \neg x_i\lor\neg x_j
  \qquad\forall i<j.
\]
Phương pháp này không dùng biến phụ, có \(\binom n2\) mệnh đề nhị phân và lan truyền trực
tiếp. Đây là lựa chọn tốt khi miền nhỏ.

\subsection{Biểu diễn tuần tự}

Đặt \(s_i=1\) khi ít nhất một trong \(x_1,\ldots,x_i\) đúng. Một biến thể
biểu diễn AMO tuần tự dùng:
\[
  \neg x_i\lor s_i,
  \qquad
  \neg s_{i-1}\lor s_i,
  \qquad
  \neg x_i\lor\neg s_{i-1}.
\]
Mệnh đề cuối cấm \(x_i\) khi tiền tố trước đó đã có một giá trị đúng. Kích thước
tuyến tính theo \(n\) \parencite{sinz2005cardinality}.

\subsection{Biểu diễn tích và biểu diễn chỉ huy}

\emph{Product encoding} (biểu diễn tích) đặt biến vào ma trận rồi dùng AMO theo
hàng và cột. \emph{Commander encoding} (biểu diễn chỉ huy) chia biến thành nhóm,
tạo một biến đại diện cho mỗi nhóm rồi biểu diễn AMO trên các đại diện. Hai cách
này tạo tính cục bộ theo khối và phù hợp với miền vừa hoặc lớn
\parencite{chen2010amo}.

\begin{center}
\captionof{table}{So sánh các bộ biểu diễn AMO cơ bản.}
\label{tab:amo-encoders}
\begin{tabularx}{\textwidth}{@{}>{\bfseries\sffamily}p{.19\textwidth}YYY@{}}
\toprule
Phép biểu diễn AMO & Biến phụ & Mệnh đề & Điểm mạnh\\
\midrule
Từng cặp & \(0\) & \(O(n^2)\) & Nhị phân, đơn giản\\
Tuần tự & \(O(n)\) & \(O(n)\) & Lan truyền theo tiền tố\\
Tích & \(O(\sqrt n)\) & \(O(n)\) & Cân bằng hai chiều\\
Chỉ huy & \(O(n/g)\) & Phụ thuộc nhóm \(g\) & Cục bộ, dễ kết hợp\\
\bottomrule
\end{tabularx}
\end{center}

\section{Sequential Counter cho AMK}

Đặt
\[
  s_{i,j}=1
  \iff
  \sum_{q=1}^{i}x_q\ge j,
  \qquad 1\le j\le k.
\]
Bộ đếm cập nhật theo hệ thức truy hồi
\[
  s_{i,j}
  \leftrightarrow
  s_{i-1,j}\lor(x_i\land s_{i-1,j-1}),
\]
với \(s_{i,0}=\top\). Cận AMK cấm trạng thái tràn:
\[
  \neg x_i\lor\neg s_{i-1,k}.
\]
Số thanh ghi và mệnh đề là \(O(nk)\). Với \(k\) nhỏ, bộ đếm tạo các quan hệ
kéo theo cục bộ và cho phép đọc nhiều ngưỡng tiền tố
\parencite{sinz2005cardinality}.

\begin{proposition}[Ngữ nghĩa của bộ đếm]
Nếu các hệ thức truy hồi được biểu diễn hai chiều thì, với mọi \(i,j\),
\[
  s_{i,j}=1\iff\sum_{q=1}^{i}x_q\ge j.
\]
\end{proposition}

\begin{proof}
Quy nạp theo \(i\). Tiền tố \(i\) đạt \(j\) khi tiền tố \(i-1\) đã đạt \(j\),
hoặc \(x_i=1\) và tiền tố trước đạt \(j-1\). Đây chính là hệ thức truy hồi.
\end{proof}

\begin{figure}[tb]
\centering
\begin{tikzpicture}[satfig,x=16mm,y=10mm]
  \foreach \i in {1,...,5}
    \node[satlabel] at (\i,2.7) {\(x_\i\)};
  \foreach \i in {1,...,5}{
    \foreach \j in {1,2}{
      \node[satstate] (s\i\j) at (\i,2-\j) {\(\scriptstyle s_{\i,\j}\)};
    }
  }
  \foreach \i in {2,...,5}{
    \draw[satedge] (s\the\numexpr\i-1\relax1)--(s\i1);
    \draw[satedge] (s\the\numexpr\i-1\relax2)--(s\i2);
  }
  \foreach \i in {2,...,5}
    \draw[satdash] (\i,2.35)--(s\i1);
  \foreach \i in {2,...,5}
    \draw[satdash] (\i,2.35) to[bend right=18] (s\i2);
  \node[satwarn,rounded corners=2pt,right=10mm of s52]
    {\(\neg s_{5,3}\)\\(trạng thái tràn bị cắt)};
\end{tikzpicture}
\caption{Lưới trạng thái của Sequential Counter cho \(n=5,k=2\). Mũi tên liền
truyền ngưỡng; mũi tên đứt đưa đầu vào mới vào hệ thức truy hồi.}
\label{fig:seq-counter-grid}
\end{figure}

\begin{workedexample}{AMK với năm đầu vào và cận hai}
Xét \(x_1+\cdots+x_5\le2\). Sau khi \(x_1=x_3=1\), bộ đếm suy ra
\[
  s_{3,2}=1.
\]
Mệnh đề chống tràn tại \(i=4\),
\[
  \neg x_4\lor\neg s_{3,2},
\]
trở thành mệnh đề đơn vị và suy ra \(\neg x_4\). Tương tự, tính đơn điệu của
tiền tố đưa \(s_{3,2}\) đến \(s_{4,2}\), rồi mệnh đề chống tràn tại \(i=5\)
suy ra \(\neg x_5\). Bộ đếm không cần đợi đầu vào đúng thứ ba mới báo xung đột; cấu trúc sẽ
loại mọi vị trí còn lại ngay khi cận đã đầy.
\end{workedexample}

\section{Totalizer}

Totalizer xây dựng một cây nhị phân. Mỗi nút trả về một vectơ đơn phân
\[
  o_j=1\iff \text{tổng trong cây con}\ge j.
\]
Khi ghép hai nút con có đầu ra \(a,b\), nút cha \(r\) thỏa
\[
  a_p\land b_q\rightarrow r_{p+q}.
\]
Chiều ngược bảo đảm đặc tả đầy đủ nếu các đầu ra được dùng bên ngoài. Cận
\(\sum x_i\le k\) chỉ cần \(\neg o_{k+1}\).

Totalizer là lựa chọn tự nhiên khi:
\begin{itemize}
  \item cần nhiều ngưỡng trên cùng một tập biến;
  \item cận được thay đổi tăng dần;
  \item cần ràng buộc đúng \(K\) bằng \(o_k\land\neg o_{k+1}\).
\end{itemize}
Totalizer cân bằng tạo tính cục bộ theo cây; kích thước có thể được cắt ở
\(k+1\) thay vì giữ toàn bộ tổng \parencite{bailleux2003cardinality}.

\begin{figure}[tb]
\centering
\begin{tikzpicture}[satfig,level distance=13mm,
  level 1/.style={sibling distance=45mm},
  level 2/.style={sibling distance=22mm}]
  \node[satstate] (r) {\(r_{1:4}\)}
    child {node[satstate] (a) {\(a_{1:2}\)}
      child {node[satbox] {\(x_1\)}}
      child {node[satbox] {\(x_2\)}}}
    child {node[satstate] (b) {\(b_{1:2}\)}
      child {node[satbox] {\(x_3\)}}
      child {node[satbox] {\(x_4\)}}};
  \node[satlabel,right=11mm of r] {\(r_j=1\iff\sum_i x_i\ge j\)};
\end{tikzpicture}
\caption{Totalizer cân bằng: mỗi nút trả về các ngưỡng của tổng trong cây con.}
\label{fig:totalizer-tree}
\end{figure}

\section{Mạng đếm}

Sorting Network sắp các đầu vào Boolean thành các đầu ra đơn phân:
\[
  o_1\ge o_2\ge\cdots\ge o_n,
  \qquad
  o_j=1\iff\sum_i x_i\ge j.
\]
Bộ so sánh
\[
  (a,b)\mapsto(a\lor b,\ a\land b)
\]
được chuyển đổi theo phương pháp Tseitin. Cận AMK là \(\neg o_{k+1}\). Mạng có cấu trúc đều, độ sâu
đa lôgarit và thuận lợi cho lan truyền song song
\parencite{batcher1968sorting,asin2011networks}.

So với Sequential Counter, mạng dùng nhiều mệnh đề hơn khi \(k\) nhỏ nhưng có độ
sâu suy diễn thấp hơn và các đầu ra toàn cục rõ ràng.

\begin{figure}[tb]
\centering
\begin{tikzpicture}[satfig,x=17mm,y=9mm]
  \node[satbox] (a) at (0,1) {\(a\)};
  \node[satbox] (b) at (0,0) {\(b\)};
  \node[satstate] (c) at (2,1) {\(\lor\)};
  \node[sataux] (d) at (2,0) {\(\land\)};
  \node[satbox] (hi) at (4,1) {\(\max(a,b)\)};
  \node[satbox] (lo) at (4,0) {\(\min(a,b)\)};
  \draw[satedge] (a)--(c);
  \draw[satedge] (b)--(c);
  \draw[satedge] (a)--(d);
  \draw[satedge] (b)--(d);
  \draw[satedge] (c)--(hi);
  \draw[satedge] (d)--(lo);
\end{tikzpicture}
\caption{Bộ so sánh Boolean là thành phần cơ sở của Sorting Network và mạng đếm.}
\label{fig:boolean-comparator}
\end{figure}

\section{Biểu diễn bằng BDD}

\emph{Binary decision diagram} (biểu đồ quyết định nhị phân, BDD) biểu diễn
trạng thái tổng tích lũy. Nút ở mức \(i\) và trạng thái \(s\) phân nhánh theo
\(x_i\). Các trạng thái vượt \(k\) đi vào nút kết thúc sai. Một biến phụ cho
mỗi nút BDD cùng các mệnh đề liên kết theo hai cung tạo thành phép biểu diễn
\parencite{bryant1986bdd,abio2012bdd}.

Ưu điểm của BDD:
\begin{itemize}
  \item xử lý tự nhiên các ràng buộc giả Boolean có trọng số;
  \item có thể đạt sức lan truyền mạnh;
  \item hợp nhất các trạng thái tương đương.
\end{itemize}
Nhược điểm là kích thước phụ thuộc thứ tự biến và số trạng thái tổng khác nhau.

\section{Ràng buộc đúng \(K\) và đầu ra chia sẻ}

EXK có thể viết thành AMK và ALK, nhưng dựng hai bộ đếm độc lập sẽ lặp trạng
thái. Nếu bộ biểu diễn trả về các đầu ra đơn phân \(o_j\), chỉ cần
\[
  o_k\land\neg o_{k+1}.
\]
Sequential Counter cũng có thể chia sẻ thanh ghi: biên \(s_{n,k}\) bảo đảm ít
nhất \(k\), còn các mệnh đề chống tràn bảo đảm nhiều nhất \(k\).

\begin{keyidea}{Định nghĩa một lần, đọc nhiều cận}
Tách \texttt{CounterDefinition} khỏi \texttt{Bound}. Cùng một cấu trúc đầu ra
có thể phục vụ AMK, ALK, EXK và vòng giải tăng dần. Đây là bước đầu của tư duy
mô-đun được phát triển ở Phần II.
\end{keyidea}

\section{Pseudo-Boolean có trọng số}

Ràng buộc tổng quát
\[
  \sum_{i=1}^{n}a_ix_i\le B,
  \qquad a_i\in\mathbb Z_{\ge0},
\]
không còn là cardinality khi trọng số khác nhau. Các hướng chính gồm:
\begin{itemize}
  \item biểu diễn bằng bộ cộng hoặc mạch cho biểu diễn nhị phân của tổng;
  \item BDD cho các trạng thái tổng có thể đạt;
  \item Totalizer tổng quát gom và cắt các tổng;
  \item phân rã bằng Sorting Network hoặc mạng đếm sau khi tách trọng số;
  \item thư viện kết hợp như PBLib chọn bộ biểu diễn theo cấu trúc
        \parencite{een2006pb,philipp2015pblib}.
\end{itemize}

Nếu hàm mục tiêu thay cận nhiều lần, các đầu ra hoặc giao diện kích hoạt quan
trọng hơn vài phần trăm kích thước ở một cận duy nhất.

\subsection{Totalizer tổng quát}

\emph{Generalized totalizer} (Totalizer tổng quát) thay vectơ đơn phân liên tiếp
bằng tập các tổng có thể đạt tại mỗi nút. Khi ghép hai nút có tập tổng \(A,B\),
nút cha giữ các giá trị
\[
  \min(a+b,B+1),\qquad a\in A,\ b\in B,
\]
trong đó \(B+1\) gom mọi trạng thái tràn. Phép biểu diễn đặc biệt hiệu quả khi
trọng số có cấu trúc và số tổng phân biệt sau phép cắt nhỏ hơn tổng số học thô
\parencite{joshi2015gte}. Hình dạng cây không còn là chi tiết vô hại: hai cách
ghép các lá có thể tạo số trạng thái trung gian rất khác nhau.

\subsection{Khai thác các nhóm AMO trong PB}

Trong nhiều mô hình, các hạng tử của ràng buộc PB đến từ một biến miền hữu hạn:
trong mỗi nhóm, nhiều nhất một literal có thể đúng. Nếu trải phẳng mọi hạng tử
như các biến độc lập, bộ biểu diễn sẽ bỏ mất thông tin này. Các phép biểu diễn
PB(AMO) đưa nhóm AMO trực tiếp vào \emph{multi-valued decision diagram}
(biểu đồ quyết định đa trị, MDD), Totalizer
tổng quát, Sequential Counter có trọng số hoặc Totalizer theo môđun
\parencite{bofill2022pbamo}. Hai lợi ích có thể xuất hiện:
\begin{itemize}
  \item ít trạng thái tổng hơn, vì tổ hợp chọn hai hạng tử cùng nhóm không bao
        giờ đạt được;
  \item lan truyền dùng được cả cận PB lẫn tính loại trừ lẫn nhau trong nhóm.
\end{itemize}

Đây là một ví dụ tổng quát của \emph{structure-aware encoding}
(phép biểu diễn có xét cấu trúc): trước khi chọn bộ biểu diễn cho một biểu thức
tuyến tính, cần giữ nguồn gốc cấu trúc của các hạng tử thay vì chỉ truyền vào
mảng hệ số.

\section{Từ lựa chọn thủ công đến lựa chọn theo bài toán thử}

Quy tắc ``\(k\) nhỏ thì dùng Sequential Counter, nhiều cận thì dùng Totalizer'' là
một đối chứng hữu ích, nhưng chưa bao quát tương tác giữa hệ số, nhóm AMO, độ
rộng miền và bộ giải. Nghiên cứu về \emph{algorithm selection}
(lựa chọn thuật toán) đã dùng các đặc trưng của ràng buộc và bài toán
thử để chọn giữa nhiều phép biểu diễn SAT, kể cả trên những lớp bài toán chưa xuất
hiện trong dữ liệu huấn luyện \parencite{ulrich2022selection}. Kết quả này thay
đổi cách xây dựng thư viện bộ biểu diễn: mỗi mô-đun cần công bố đủ siêu dữ liệu để
một bộ chọn ra quyết định.

\begin{center}
\captionof{table}{Nhóm đặc trưng để lựa chọn bộ biểu diễn theo bài toán thử.}
\label{tab:selector-features}
\begin{tabularx}{\textwidth}{@{}>{\bfseries\sffamily}p{.22\textwidth}YY@{}}
\toprule
Nhóm đặc trưng & Ví dụ & Tác động tiềm năng\\
\midrule
Kích thước & \(n,k,B\), số hạng tử & Loại bộ biểu diễn chắc chắn quá lớn\\
Trọng số & Giá trị lớn nhất, ước chung lớn nhất, số tổng phân biệt & Không gian trạng thái BDD/GT\\
Cấu trúc & Nhóm AMO, hệ số lặp & Khả năng nén trạng thái\\
Giao diện & Số cận, thêm đầu vào, hiện thực hóa (reification) & Giải tăng dần và kết hợp mô-đun\\
Mô hình & Mật độ, đối xứng, đồ thị ràng buộc & Tương tác với học mệnh đề\\
\bottomrule
\end{tabularx}
\end{center}

Bộ chọn không loại bỏ nhu cầu chứng minh tính đúng: công cụ chỉ lựa chọn giữa các
mô-đun đã được kiểm chứng. Việc đánh giá cũng phải tách các họ dùng để huấn
luyện khỏi các họ dùng để kiểm thử; nếu không, bộ chọn có thể chỉ nhận dạng tên
bộ kiểm chuẩn thay vì học cấu trúc của phép biểu diễn.

\section{Bản đồ lựa chọn}

\begin{center}
\captionof{table}{Gợi ý lựa chọn bộ biểu diễn ràng buộc đếm/PB.}
\label{tab:cardinality-choice}
\begin{tabularx}{\textwidth}{@{}>{\bfseries\sffamily}p{.26\textwidth}YY@{}}
\toprule
Tình huống & Lựa chọn khởi đầu & Lý do\\
\midrule
AMO, \(n\) nhỏ & Từng cặp & Không biến phụ, mệnh đề nhị phân\\
AMO/AMK, \(k\) nhỏ & Tuần tự & \(O(nk)\), lan truyền theo tiền tố\\
Nhiều ngưỡng & Totalizer/mạng đếm & Tái sử dụng đầu ra\\
Đúng \(K\) & Một vectơ đầu ra chung & Không lặp AMK/ALK\\
PB có trọng số, ít trạng thái & BDD & Hợp nhất trạng thái tổng\\
PB có trọng số tổng quát & PBLib/kết hợp & Chọn theo bài toán thử\\
PB có nhóm AMO & Bộ biểu diễn PB(AMO) & Không hiện thực hóa (reification) trạng thái bất khả\\
Cận thay đổi tăng dần & Totalizer/đầu ra thứ tự & Giao diện cận ổn định\\
\bottomrule
\end{tabularx}
\end{center}

Đây là cấu hình khởi đầu; Chương 4 trình bày cách đánh giá bộ biểu diễn trên phân
phối bài toán mục tiêu; Chương 5--6 phát triển các bộ đếm khi nhiều
ràng buộc chia sẻ cấu trúc.

\section{Bài học thiết kế}

\begin{enumerate}
  \item Chọn bộ biểu diễn theo cận và đầu ra cần đọc, không chỉ theo tên ràng buộc.
  \item Biểu diễn AMO từng cặp vẫn là đối chứng mạnh khi miền nhỏ.
  \item Sequential Counter phù hợp với ngưỡng nhỏ và cấu trúc tiền tố.
  \item Totalizer/mạng đếm thuận lợi khi cần nhiều cận hoặc giải tăng dần.
  \item Ràng buộc đúng \(K\) nên chia sẻ một cấu trúc tổng nếu có thể.
\end{enumerate}

\subsection*{Câu hỏi nghiên cứu}
\begin{enumerate}
  \item Dựng Sequential Counter cho \(n=5,k=2\) và liệt kê các thanh ghi.
  \item So sánh EXK bằng hai bộ đếm độc lập với một Totalizer chung.
  \item Với trọng số \((1,2,3,5)\), vẽ các trạng thái BDD cho cận \(B=5\).
\end{enumerate}
